% 模板 9：圆形编号 + 卡片分层（高级经管/实证风格）
% 统一精致配色：蓝色主节点+teal子节点+微阴影
% 编译命令：xelatex demo_roadmap_research_premium.tex
\documentclass[border=5pt]{standalone}
\usepackage{ctex}
\usepackage{tikz}
\usetikzlibrary{arrows.meta, positioning, calc, shadows}

\begin{document}
\begin{tikzpicture}[
    phasenum/.style={circle, fill=blue!70, minimum size=20pt,
        font=\scriptsize\bfseries, text=white, inner sep=0pt,
        drop shadow={opacity=0.15, shadow xshift=0.5pt, shadow yshift=-0.5pt}},
    phasename/.style={font=\small\bfseries, color=blue!70, anchor=west},
    method/.style={rectangle, rounded corners=3pt,
        minimum width=3.4cm, minimum height=0.7cm, align=center,
        draw=blue!80, line width=0.7pt, fill=blue!6, font=\small,
        drop shadow={opacity=0.15, shadow xshift=0.5pt, shadow yshift=-0.5pt}},
    tool/.style={rectangle, rounded corners=2pt,
        minimum width=1.8cm, minimum height=0.5cm, align=center,
        draw=teal!70, line width=0.5pt, fill=white, font=\scriptsize},
    outputtag/.style={fill=blue!65, rounded corners=10pt,
        minimum width=1.5cm, minimum height=0.38cm, align=center,
        font=\tiny\bfseries, text=white, inner sep=2pt},
    pipe/.style={-stealth, line width=1.4pt, color=blue!70},
    inner/.style={-stealth, line width=0.5pt, color=gray!50},
    card/.style={fill=gray!2, rounded corners=4pt,
        draw=gray!25, dashed, line width=0.7pt},
]

% ====== Phase 1: 研究设计 ======
\node[phasenum] at (0, 2.6) {1};
\node[phasename] at (0.6, 2.6) {研究设计};
\fill[card] (-0.5, 2.2) rectangle (15.5, -0.5);

\node[method] (rq) at (3.2, 1.2) {研究问题提出};
\node[method] (lit) at (7.2, 1.2) {系统文献综述};
\node[method] (hypo) at (11.2, 1.2) {假设与框架构建};
\draw[inner] (rq) -- (lit); \draw[inner] (lit) -- (hypo);

\node[tool] at (3.2, 0.2) {文献计量};
\node[tool] at (5.5, 0.2) {知识图谱};
\node[tool] at (8.2, 0.2) {理论推演};
\node[tool] at (11.2, 0.2) {概念模型};
\node[outputtag] at (14.2, 1.2) {理论模型};

\draw[pipe] (7.2, -0.5) -- (7.2, -1.4);

% ====== Phase 2: 数据准备 ======
\node[phasenum] at (0, -1.5) {2};
\node[phasename] at (0.6, -1.5) {数据准备};
\fill[card] (-0.5, -1.9) rectangle (15.5, -6.2);

\node[method] (collect) at (3.2, -2.5) {数据采集};
\node[method] (clean) at (7.2, -2.5) {数据清洗与融合};
\node[method] (eda) at (11.2, -2.5) {探索性分析};
\draw[inner] (collect) -- (clean); \draw[inner] (clean) -- (eda);

\node[tool] at (2, -3.5) {多源采集};
\node[tool] at (4.2, -3.5) {缺失值处理};
\node[tool] at (6.5, -3.5) {异常值检测};
\node[tool] at (9, -3.5) {分布检验};
\node[tool] at (11.2, -3.5) {可视化};

\node[method] (vardef) at (3.2, -4.6) {变量操作化};
\node[method] (feateng) at (7.2, -4.6) {变量构建};
\node[method] (descstat) at (11.2, -4.6) {描述性统计};
\draw[inner] (vardef) -- (feateng); \draw[inner] (feateng) -- (descstat);

\node[tool] at (2.2, -5.5) {标准化};
\node[tool] at (4.5, -5.5) {交互项};
\node[tool] at (7.2, -5.5) {多重共线性};
\node[tool] at (10, -5.5) {相关矩阵};
\node[outputtag] at (14.2, -2.5) {分析数据集};

\draw[pipe] (7.2, -6.2) -- (7.2, -7.1);

% ====== Phase 3: 实证分析 ======
\node[phasenum] at (0, -7.2) {3};
\node[phasename] at (0.6, -7.2) {实证分析};
\fill[card] (-0.5, -7.6) rectangle (15.5, -15.5);

\node[method] (spec) at (3.2, -8.2) {模型设定};
\node[method] (ident) at (7.2, -8.2) {识别策略};
\node[method] (est) at (11.2, -8.2) {参数估计};
\draw[inner] (spec) -- (ident); \draw[inner] (ident) -- (est);

\node[tool] at (1.8, -9.2) {固定效应};
\node[tool] at (4, -9.2) {DID};
\node[tool] at (6.2, -9.2) {工具变量};
\node[tool] at (8.4, -9.2) {PSM};
\node[tool] at (10.4, -9.2) {RDD};
\node[tool] at (12.4, -9.2) {GMM};

\node[method] (mech2) at (3.2, -10.5) {机制检验};
\node[method] (hetero2) at (7.2, -10.5) {异质性分析};
\node[method] (ext) at (11.2, -10.5) {拓展分析};
\draw[inner] (spec.south) -- ++(0,-0.5) -- (mech2.north);
\draw[inner] (ident.south) -- ++(0,-0.5) -- (hetero2.north);
\draw[inner] (est.south) -- ++(0,-0.5) -- (ext.north);

\node[tool] at (1.8, -11.5) {中介效应};
\node[tool] at (4, -11.5) {调节效应};
\node[tool] at (6.5, -11.5) {分组回归};
\node[tool] at (9, -11.5) {分位数};
\node[tool] at (11.2, -11.5) {空间溢出};
\node[tool] at (13.2, -11.5) {动态面板};

\node[method] (robust2) at (7.2, -12.8) {稳健性检验};
\draw[inner] (mech2.south) -- ++(0,-0.3) -| ([xshift=-8pt]robust2.north);
\draw[inner] (ext.south) -- ++(0,-0.3) -| ([xshift=8pt]robust2.north);

\node[tool] at (2.2, -13.8) {替换变量};
\node[tool] at (4.8, -13.8) {缩尾处理};
\node[tool] at (7.5, -13.8) {安慰剂检验};
\node[tool] at (10.2, -13.8) {子样本};
\node[tool] at (12.6, -13.8) {Heckman};
\node[outputtag] at (14.2, -8.2) {估计结果};

\draw[pipe] (7.2, -15.5) -- (7.2, -16.4);

% ====== Phase 4: 结论建议 ======
\node[phasenum] at (0, -16.5) {4};
\node[phasename] at (0.6, -16.5) {结论建议};
\fill[card] (-0.5, -16.9) rectangle (15.5, -20.0);

\node[method] (findings) at (3.2, -17.5) {核心发现};
\node[method] (contrib) at (7.2, -17.5) {理论贡献};
\node[method] (policy2) at (11.2, -17.5) {政策启示};
\draw[inner] (findings) -- (contrib); \draw[inner] (contrib) -- (policy2);

\node[method] (limit) at (5, -18.8) {研究局限};
\node[method] (future2) at (9.5, -18.8) {未来方向};
\draw[inner] (findings.south) -- ++(0,-0.3) -| (limit.north);
\draw[inner] (policy2.south) -- ++(0,-0.3) -| (future2.north);
\node[outputtag] at (14.2, -17.5) {研究成果};

\end{tikzpicture}
\end{document}
